\DocumentMetadata{
  pdfversion=2.0,
  pdfstandard=ua-2,
  lang=en-US,
  tagging=on,
  tagging-setup={math/setup=mathml-SE,tabsorder=structure}
}
\documentclass[11pt,letterpaper]{article}
\usepackage[margin=0.68in,includefoot]{geometry}
\usepackage{newcomputermodern}
\usepackage{amsmath,amscd,mathtools}
\usepackage{tikz,adjustbox}
\providecommand{\square}{\mdlgwhtsquare}
\usepackage[hidelinks]{hyperref}
\hypersetup{
  pdftitle={Complex Analysis PhD Exam, August 2023},
  pdfauthor={University of Florida Department of Mathematics},
  pdfsubject={Graduate mathematics examination},
  pdfdisplaydoctitle=true,
  bookmarksopen=true
}
\setcounter{secnumdepth}{0}
\setlength{\parindent}{0pt}
\setlength{\parskip}{0.28em}
\setlength{\emergencystretch}{3em}
\setlength{\leftmargini}{2.15em}
\setlength{\leftmarginii}{2.65em}
\setlength{\leftmarginiii}{3em}
\pagestyle{empty}
\raggedbottom
\allowdisplaybreaks
\NewTaggingSocketPlug{block/list/label}{examproblem}{%
  \tagstructbegin{tag=Lbl}\tagstructend%
  \tagstructbegin{tag=\LBody}%
  \tagstructbegin{tag=\ProblemHeadingTag,title-o={\ProblemHeadingText}}%
  \tagmcbegin{tag=\ProblemHeadingTag,actualtext={\ProblemHeadingText}}%
  #2%
  \tagmcend%
  \tagstructend%
}
\newenvironment{examproblems}[2]{%
  \def\ProblemHeadingTag{#2}%
  \AssignTaggingSocketPlug{block/list/label}{examproblem}%
  \begin{enumerate}%
  \setcounter{enumi}{#1}%
  \renewcommand{\labelenumi}{\textbf{\theenumi.}}%
  \setlength{\topsep}{0.35em}%
  \setlength{\partopsep}{0pt}%
  \setlength{\itemsep}{0.48em}%
  \setlength{\parsep}{0.18em}%
}{\end{enumerate}}
\newenvironment{examsubparts}{%
  \AssignTaggingSocketPlug{block/list/label}{default}%
  \begin{enumerate}%
  \setlength{\topsep}{0.18em}%
  \setlength{\partopsep}{0pt}%
  \setlength{\itemsep}{0.16em}%
  \setlength{\parsep}{0.1em}%
}{\end{enumerate}}
\newenvironment{examinstructions}{%
  \par\begingroup\itshape
}{%
  \par\endgroup\vspace{0.18em}
}
\makeatletter
\renewcommand\subsection{\@startsection{subsection}{2}{0pt}%
  {-1.25ex plus -.25ex minus -.1ex}%
  {0.55ex plus .1ex}%
  {\normalfont\large\bfseries}}
\renewcommand{\@maketitle}{%
  \newpage\null\vskip -1em%
  {\footnotesize\bfseries
    \href{https://www.ufl.edu/}{University of Florida}, \href{https://math.ufl.edu/}{Department of Mathematics}\par}%
  \vskip -0.5em%
  \tagstructbegin{tag=H1,title={Complex Analysis PhD Exam, August 2023}}%
  \tagpdfparaOff%
  \tagmcbegin{tag=H1}%
    {\LARGE\bfseries\href{https://gma.math.ufl.edu/exams/complex-analysis-phd/}{Complex Analysis PhD Exam}, August 2023\par}%
  \tagmcend%
  \tagpdfparaOn%
  \tagstructend%
  \vskip 0.25em%
  \tagmcbegin{artifact}%
    \hrule height 1pt%
  \tagmcend%
  \vskip 1em%
}
\makeatother
\title{Complex Analysis PhD Exam, August 2023}
\author{}
\date{}
\begin{document}
\maketitle
\thispagestyle{empty}
\begin{examinstructions}
Answer FIVE (out of eight) questions in detail, explaining your work in a precise and logical way. State carefully any results used without proof.
\end{examinstructions}
\begin{examproblems}{0}{H2}
\def\ProblemHeadingText{Problem 1.}
\item
Let \(S=\{z\in\mathbb{C}:|z|<2,\ |z-1|>1\}\) and \(\mathbb{D}=\{z\in\mathbb{C}:|z|<1\}\). Find an explicit analytic bijection \(f:\mathbb{D}\to S\).
\def\ProblemHeadingText{Problem 2.}
\item
Fix \(\lambda>1\). Show that the equation \(\lambda-z-e^{-z}=0\) has exactly one solution in the right half plane, \(\mathbb{H}=\{z\in\mathbb{C}:\operatorname{real} z>0\}\).
\def\ProblemHeadingText{Problem 3.}
\item
Let~\(A\) be an \(n\times n\) matrix and \(p(z)=\det(zI_n-A)\). Show \(p(A)=0\), where \(I_n\) is the \(n\times n\) identity matrix.
\def\ProblemHeadingText{Problem 4.}
\item
Show \(e^z-z^2=0\) has infinitely many solutions.
\def\ProblemHeadingText{Problem 5.}
\item
Find all solutions to \(f^n+g^n=1\) for positive integers \(n\geq 2\) and entire functions~\(f\) and~\(g\).
\def\ProblemHeadingText{Problem 6.}
\item
Show \(\displaystyle\prod_{n=2}^{\infty}\left(1-\frac{1}{n^2}\right)=\frac{1}{2}\).
\def\ProblemHeadingText{Problem 7.}
\item
Suppose \(\Omega\subseteq\mathbb{C}\) is a region (non-empty, open and connected), \((f_n)\) is a locally bounded sequence of analytic functions \(f_n:\Omega\to\mathbb{C}\) and \(S\subseteq\Omega\) has an accumulation point in~\(\Omega\). Show, if \(f:\Omega\to\mathbb{C}\) is an analytic function and \((f_n(s))\) converges to \(f(s)\) for each \(s\in S\), then \((f_n)\) converges to~\(f\) in \(H(\Omega)\) (uniformly on compact sets).
\def\ProblemHeadingText{Problem 8.}
\item
Let~\(\zeta\) denote the Riemann zeta function and~\(\sigma\) the divisor function. Thus, for positive integers, \(\sigma(k)\) is the number of divisors of~\(k\). Show, if real \(z>1\), then
\[
\zeta(z)^2=\sum_{k=1}^{\infty}\frac{\sigma(k)}{k^z}.
\]
\end{examproblems}
\end{document}
